\documentclass[10pt,a4paper]{amsart}
\usepackage{amsmath,amssymb,amsthm,graphicx}
% For commutative diagrams we use Xy-pic
\usepackage[all]{xy}

% We use 2cell for 2-commutative diagrams.
\xyoption{2cell}
% Theorem environments.
%
\theoremstyle{plain}
\newtheorem{theorem}[subsection]{Theorem}
\newtheorem{proposition}[subsection]{Proposition}
\newtheorem{lemma}[subsection]{Lemma}

\theoremstyle{definition}
\newtheorem{definition}[subsection]{Definition}
\newtheorem{example}[subsection]{Example}
\newtheorem{exercise}[subsection]{Exercise}
\newtheorem{situation}[subsection]{Situation}

\theoremstyle{remark}
\newtheorem{remark}[subsection]{Remark}
\newtheorem{remarks}[subsection]{Remarks}

\numberwithin{equation}{subsection}

% Macros
%
\def\lim{\mathop{\mathrm{lim}}\nolimits}
\def\colim{\mathop{\mathrm{colim}}\nolimits}
\def\Spec{\mathop{\mathrm{Spec}}}
\def\Hom{\mathop{\mathrm{Hom}}\nolimits}
\def\Ext{\mathop{\mathrm{Ext}}\nolimits}
\def\SheafHom{\mathop{\mathcal{H}\!\mathit{om}}\nolimits}
\def\SheafExt{\mathop{\mathcal{E}\!\mathit{xt}}\nolimits}
\def\Sch{\mathit{Sch}}
\def\Mor{\mathop{\mathrm{Mor}}\nolimits}
\def\Ob{\mathop{\mathrm{Ob}}\nolimits}
\def\Sh{\mathop{\mathit{Sh}}\nolimits}
\def\NL{\mathop{N\!L}\nolimits}
\def\CH{\mathop{\mathrm{CH}}\nolimits}
\def\proetale{{pro\text{-}\acute{e}tale}}
\def\etale{{\acute{e}tale}}
\def\QCoh{\mathit{QCoh}}
\def\Ker{\mathop{\mathrm{Ker}}}
\def\Im{\mathop{\mathrm{Im}}}
\def\Coker{\mathop{\mathrm{Coker}}}
\def\Coim{\mathop{\mathrm{Coim}}}

% Boxtimes
%
\DeclareMathSymbol{\boxtimes}{\mathbin}{AMSa}{"02}

%
% Macros for moduli stacks/spaces
%
\def\QCohstack{\mathcal{QC}\!\mathit{oh}}
\def\Cohstack{\mathcal{C}\!\mathit{oh}}
\def\Spacesstack{\mathcal{S}\!\mathit{paces}}
\def\Quotfunctor{\mathrm{Quot}}
\def\Hilbfunctor{\mathrm{Hilb}}
\def\Curvesstack{\mathcal{C}\!\mathit{urves}}
\def\Polarizedstack{\mathcal{P}\!\mathit{olarized}}
\def\Complexesstack{\mathcal{C}\!\mathit{omplexes}}
% \Pic is the operator that assigns to X its picard group, usage \Pic(X)
% \Picardstack_{X/B} denotes the Picard stack of X over B
% \Picardfunctor_{X/B} denotes the Picard functor of X over B
\def\Pic{\mathop{\mathrm{Pic}}\nolimits}
\def\Picardstack{\mathcal{P}\!\mathit{ic}}
\def\Picardfunctor{\mathrm{Pic}}
\def\Deformationcategory{\mathcal{D}\!\mathit{ef}}
\newenvironment{reference}{}{}
\newenvironment{slogan}{}{}
\newenvironment{history}{}{}
\usepackage{cleveref}
\crefname{originallemma}{original lemma}{original lemmas}
\Crefname{originallemma}{Original lemma}{Original lemmas}
\makeatletter
\newlabel{dualizing-lemma-local-cohomology-noetherian}{{0955}{source}{}{sourcepointer-1}{}}
\newlabel{dualizing-lemma-local-cohomology-noetherian@cref}{{[originallemma][0][]0955}{[1][0][]source}{}{}{}}
\makeatother
\crefname{originallemma}{original lemma}{original lemmas}
\Crefname{originallemma}{Original lemma}{Original lemmas}
\makeatletter
\newlabel{dualizing-lemma-complete-and-local}{{0A6W}{source}{}{sourcepointer-2}{}}
\newlabel{dualizing-lemma-complete-and-local@cref}{{[originallemma][0][]0A6W}{[1][0][]source}{}{}{}}
\makeatother
\crefname{originallemma}{original lemma}{original lemmas}
\Crefname{originallemma}{Original lemma}{Original lemmas}
\makeatletter
\newlabel{dualizing-lemma-local-cohomology-ss}{{0BJC}{source}{}{sourcepointer-3}{}}
\newlabel{dualizing-lemma-local-cohomology-ss@cref}{{[originallemma][0][]0BJC}{[1][0][]source}{}{}{}}
\makeatother
\crefname{originallemma}{original lemma}{original lemmas}
\Crefname{originallemma}{Original lemma}{Original lemmas}
\makeatletter
\newlabel{more-algebra-lemma-derived-completion}{{091V}{source}{}{sourcepointer-4}{}}
\newlabel{more-algebra-lemma-derived-completion@cref}{{[originallemma][0][]091V}{[1][0][]source}{}{}{}}
\makeatother
\crefname{originallemma}{original lemma}{original lemmas}
\Crefname{originallemma}{Original lemma}{Original lemmas}
\makeatletter
\newlabel{dualizing-lemma-compute-local-cohomology-noetherian}{{0956}{source}{}{sourcepointer-5}{}}
\newlabel{dualizing-lemma-compute-local-cohomology-noetherian@cref}{{[originallemma][0][]0956}{[1][0][]source}{}{}{}}
\makeatother
\crefname{originallemma}{original lemma}{original lemmas}
\Crefname{originallemma}{Original lemma}{Original lemmas}
\makeatletter
\newlabel{local-cohomology-lemma-kill-local-cohomology-at-prime}{{0EFE}{source}{}{sourcepointer-6}{}}
\newlabel{local-cohomology-lemma-kill-local-cohomology-at-prime@cref}{{[originallemma][0][]0EFE}{[1][0][]source}{}{}{}}
\makeatother
\crefname{originallemma}{original lemma}{original lemmas}
\Crefname{originallemma}{Original lemma}{Original lemmas}
\makeatletter
\newlabel{local-cohomology-lemma-depth-function}{{0DWY}{source}{}{sourcepointer-7}{}}
\newlabel{local-cohomology-lemma-depth-function@cref}{{[originallemma][0][]0DWY}{[1][0][]source}{}{}{}}
\makeatother
\crefname{originallemma}{original lemma}{original lemmas}
\Crefname{originallemma}{Original lemma}{Original lemmas}
\makeatletter
\newlabel{duality-lemma-sitting-in-degrees}{{0ECM}{source}{}{sourcepointer-8}{}}
\newlabel{duality-lemma-sitting-in-degrees@cref}{{[originallemma][0][]0ECM}{[1][0][]source}{}{}{}}
\makeatother
\crefname{originallemma}{original lemma}{original lemmas}
\Crefname{originallemma}{Original lemma}{Original lemmas}
\makeatletter
\newlabel{local-cohomology-lemma-isomorphism}{{0DWV}{source}{}{sourcepointer-9}{}}
\newlabel{local-cohomology-lemma-isomorphism@cref}{{[originallemma][0][]0DWV}{[1][0][]source}{}{}{}}
\makeatother
\crefname{originallemma}{original lemma}{original lemmas}
\Crefname{originallemma}{Original lemma}{Original lemmas}
\makeatletter
\newlabel{dualizing-lemma-depth}{{0AVZ}{source}{}{sourcepointer-10}{}}
\newlabel{dualizing-lemma-depth@cref}{{[originallemma][0][]0AVZ}{[1][0][]source}{}{}{}}
\makeatother
\crefname{originallemma}{original lemma}{original lemmas}
\Crefname{originallemma}{Original lemma}{Original lemmas}
\makeatletter
\newlabel{local-cohomology-lemma-adjoint-ext}{{0EF1}{source}{}{sourcepointer-11}{}}
\newlabel{local-cohomology-lemma-adjoint-ext@cref}{{[originallemma][0][]0EF1}{[1][0][]source}{}{}{}}
\makeatother
\crefname{originalproposition}{original proposition}{original propositions}
\Crefname{originalproposition}{Original proposition}{Original propositions}
\makeatletter
\newlabel{local-cohomology-proposition-annihilator}{{0EFC}{source}{}{sourcepointer-12}{}}
\newlabel{local-cohomology-proposition-annihilator@cref}{{[originalproposition][0][]0EFC}{[1][0][]source}{}{}{}}
\makeatother
\crefname{originalproposition}{original proposition}{original propositions}
\Crefname{originalproposition}{Original proposition}{Original propositions}
\makeatletter
\newlabel{local-cohomology-proposition-finiteness}{{0EFD}{source}{}{sourcepointer-13}{}}
\newlabel{local-cohomology-proposition-finiteness@cref}{{[originalproposition][0][]0EFD}{[1][0][]source}{}{}{}}
\makeatother
\crefname{originallemma}{original lemma}{original lemmas}
\Crefname{originallemma}{Original lemma}{Original lemmas}
\makeatletter
\newlabel{algebra-lemma-radical-completion}{{05GI}{source}{}{sourcepointer-14}{}}
\newlabel{algebra-lemma-radical-completion@cref}{{[originallemma][0][]05GI}{[1][0][]source}{}{}{}}
\makeatother
\crefname{originallemma}{original lemma}{original lemmas}
\Crefname{originallemma}{Original lemma}{Original lemmas}
\makeatletter
\newlabel{local-cohomology-lemma-cd-local}{{0DXB}{source}{}{sourcepointer-15}{}}
\newlabel{local-cohomology-lemma-cd-local@cref}{{[originallemma][0][]0DXB}{[1][0][]source}{}{}{}}
\makeatother
\crefname{originallemma}{original lemma}{original lemmas}
\Crefname{originallemma}{Original lemma}{Original lemmas}
\makeatletter
\newlabel{local-cohomology-lemma-cd-bound-dim-local}{{0DXF}{source}{}{sourcepointer-16}{}}
\newlabel{local-cohomology-lemma-cd-bound-dim-local@cref}{{[originallemma][0][]0DXF}{[1][0][]source}{}{}{}}
\makeatother
\crefname{originallemma}{original lemma}{original lemmas}
\Crefname{originallemma}{Original lemma}{Original lemmas}
\makeatletter
\newlabel{dualizing-lemma-completion-local}{{0EEW}{source}{}{sourcepointer-17}{}}
\newlabel{dualizing-lemma-completion-local@cref}{{[originallemma][0][]0EEW}{[1][0][]source}{}{}{}}
\makeatother
\crefname{originallemma}{original lemma}{original lemmas}
\Crefname{originallemma}{Original lemma}{Original lemmas}
\makeatletter
\newlabel{dualizing-lemma-completion-local-H0}{{0EEX}{source}{}{sourcepointer-18}{}}
\newlabel{dualizing-lemma-completion-local-H0@cref}{{[originallemma][0][]0EEX}{[1][0][]source}{}{}{}}
\makeatother
\crefname{originallemma}{original lemma}{original lemmas}
\Crefname{originallemma}{Original lemma}{Original lemmas}
\makeatletter
\newlabel{algebra-lemma-depth-dim-associated-primes}{{0BK4}{source}{}{sourcepointer-19}{}}
\newlabel{algebra-lemma-depth-dim-associated-primes@cref}{{[originallemma][0][]0BK4}{[1][0][]source}{}{}{}}
\makeatother
\crefname{originalsection}{original section}{original sections}
\Crefname{originalsection}{Original section}{Original sections}
\makeatletter
\newlabel{local-cohomology-remark-upshot}{{0EF4}{source}{}{sourcepointer-20}{}}
\newlabel{local-cohomology-remark-upshot@cref}{{[originalsection][0][]0EF4}{[1][0][]source}{}{}{}}
\makeatother
\title{Algebraization of truncated derived local cohomology under depth and dimension bounds}
\author{The Stacks Project Authors}
\date{stacks-project:algebraization.tex:local-snapshot-20261011; GNU FDL 1.2 or later}
\begin{document}
\maketitle
\noindent\textit{Editorial note.} All five hypothesis groups, including the three-way alternative in condition5, remain verbatim. The source future-reference footnote is not used as a proof dependency.\par\medskip
\noindent\textit{Editorial note.} Public310 is an annihilator/depth equivalence. This task constructs a support ideal making a derived-completion comparison an isomorphism in a specified range; it is not that annihilator theorem again.\par\medskip
\noindent\textit{Allowed standard prerequisites (editorial).}\begin{itemize}
\item Explicit prerequisite boundary. Noetherian local cohomology, dualizing-complex dimension functions, derivedI-completion, depth and support, local duality and annihilator criteria. The two local kill-completion/kill-colimit proofs are retained.
\end{itemize}\medskip
\ifdefined\hypertarget\hypertarget{sourcepointer-1}{}\fi
\noindent\textit{Source reference.} Original lemma 0955: focused-seeds/AGP/sources/stacks-snapshot-20261011/tags.txt, source lines 11573--11573. Stable Stacks tag for an expressly allowed foundational or notational input; the counted construction remains the two concrete core obligations. The retained local core is not replaced by this pointer. See standard\_prerequisites.\par\medskip
\ifdefined\hypertarget\hypertarget{sourcepointer-2}{}\fi
\noindent\textit{Source reference.} Original lemma 0A6W: focused-seeds/AGP/sources/stacks-snapshot-20261011/tags.txt, source lines 12859--12859. Stable Stacks tag for an expressly allowed foundational or notational input; the counted construction remains the two concrete core obligations. The retained local core is not replaced by this pointer. See standard\_prerequisites.\par\medskip
\ifdefined\hypertarget\hypertarget{sourcepointer-3}{}\fi
\noindent\textit{Source reference.} Original lemma 0BJC: focused-seeds/AGP/sources/stacks-snapshot-20261011/tags.txt, source lines 14520--14520. Stable Stacks tag for an expressly allowed foundational or notational input; the counted construction remains the two concrete core obligations. The retained local core is not replaced by this pointer. See standard\_prerequisites.\par\medskip
\ifdefined\hypertarget\hypertarget{sourcepointer-4}{}\fi
\noindent\textit{Source reference.} Original lemma 091V: focused-seeds/AGP/sources/stacks-snapshot-20261011/tags.txt, source lines 11458--11458. Stable Stacks tag for an expressly allowed foundational or notational input; the counted construction remains the two concrete core obligations. The retained local core is not replaced by this pointer. See standard\_prerequisites.\par\medskip
\ifdefined\hypertarget\hypertarget{sourcepointer-5}{}\fi
\noindent\textit{Source reference.} Original lemma 0956: focused-seeds/AGP/sources/stacks-snapshot-20261011/tags.txt, source lines 11574--11574. Stable Stacks tag for an expressly allowed foundational or notational input; the counted construction remains the two concrete core obligations. The retained local core is not replaced by this pointer. See standard\_prerequisites.\par\medskip
\ifdefined\hypertarget\hypertarget{sourcepointer-6}{}\fi
\noindent\textit{Source reference.} Original lemma 0EFE: focused-seeds/AGP/sources/stacks-snapshot-20261011/tags.txt, source lines 18059--18059. Stable Stacks tag for an expressly allowed foundational or notational input; the counted construction remains the two concrete core obligations. The retained local core is not replaced by this pointer. See standard\_prerequisites.\par\medskip
\ifdefined\hypertarget\hypertarget{sourcepointer-7}{}\fi
\noindent\textit{Source reference.} Original lemma 0DWY: focused-seeds/AGP/sources/stacks-snapshot-20261011/tags.txt, source lines 17413--17413. Stable Stacks tag for an expressly allowed foundational or notational input; the counted construction remains the two concrete core obligations. The retained local core is not replaced by this pointer. See standard\_prerequisites.\par\medskip
\ifdefined\hypertarget\hypertarget{sourcepointer-8}{}\fi
\noindent\textit{Source reference.} Original lemma 0ECM: focused-seeds/AGP/sources/stacks-snapshot-20261011/tags.txt, source lines 17962--17962. Stable Stacks tag for an expressly allowed foundational or notational input; the counted construction remains the two concrete core obligations. The retained local core is not replaced by this pointer. See standard\_prerequisites.\par\medskip
\ifdefined\hypertarget\hypertarget{sourcepointer-9}{}\fi
\noindent\textit{Source reference.} Original lemma 0DWV: focused-seeds/AGP/sources/stacks-snapshot-20261011/tags.txt, source lines 17410--17410. Stable Stacks tag for an expressly allowed foundational or notational input; the counted construction remains the two concrete core obligations. The retained local core is not replaced by this pointer. See standard\_prerequisites.\par\medskip
\ifdefined\hypertarget\hypertarget{sourcepointer-10}{}\fi
\noindent\textit{Source reference.} Original lemma 0AVZ: focused-seeds/AGP/sources/stacks-snapshot-20261011/tags.txt, source lines 13702--13702. Stable Stacks tag for an expressly allowed foundational or notational input; the counted construction remains the two concrete core obligations. The retained local core is not replaced by this pointer. See standard\_prerequisites.\par\medskip
\ifdefined\hypertarget\hypertarget{sourcepointer-11}{}\fi
\noindent\textit{Source reference.} Original lemma 0EF1: focused-seeds/AGP/sources/stacks-snapshot-20261011/tags.txt, source lines 18046--18046. Stable Stacks tag for an expressly allowed foundational or notational input; the counted construction remains the two concrete core obligations. The retained local core is not replaced by this pointer. See standard\_prerequisites.\par\medskip
\ifdefined\hypertarget\hypertarget{sourcepointer-12}{}\fi
\noindent\textit{Source reference.} Original proposition 0EFC: focused-seeds/AGP/sources/stacks-snapshot-20261011/tags.txt, source lines 18057--18057. Stable Stacks tag for an expressly allowed foundational or notational input; the counted construction remains the two concrete core obligations. The retained local core is not replaced by this pointer. See standard\_prerequisites.\par\medskip
\ifdefined\hypertarget\hypertarget{sourcepointer-13}{}\fi
\noindent\textit{Source reference.} Original proposition 0EFD: focused-seeds/AGP/sources/stacks-snapshot-20261011/tags.txt, source lines 18058--18058. Stable Stacks tag for an expressly allowed foundational or notational input; the counted construction remains the two concrete core obligations. The retained local core is not replaced by this pointer. See standard\_prerequisites.\par\medskip
\ifdefined\hypertarget\hypertarget{sourcepointer-14}{}\fi
\noindent\textit{Source reference.} Original lemma 05GI: focused-seeds/AGP/sources/stacks-snapshot-20261011/tags.txt, source lines 7062--7062. Stable Stacks tag for an expressly allowed foundational or notational input; the counted construction remains the two concrete core obligations. The retained local core is not replaced by this pointer. See standard\_prerequisites.\par\medskip
\ifdefined\hypertarget\hypertarget{sourcepointer-15}{}\fi
\noindent\textit{Source reference.} Original lemma 0DXB: focused-seeds/AGP/sources/stacks-snapshot-20261011/tags.txt, source lines 17426--17426. Stable Stacks tag for an expressly allowed foundational or notational input; the counted construction remains the two concrete core obligations. The retained local core is not replaced by this pointer. See standard\_prerequisites.\par\medskip
\ifdefined\hypertarget\hypertarget{sourcepointer-16}{}\fi
\noindent\textit{Source reference.} Original lemma 0DXF: focused-seeds/AGP/sources/stacks-snapshot-20261011/tags.txt, source lines 17430--17430. Stable Stacks tag for an expressly allowed foundational or notational input; the counted construction remains the two concrete core obligations. The retained local core is not replaced by this pointer. See standard\_prerequisites.\par\medskip
\ifdefined\hypertarget\hypertarget{sourcepointer-17}{}\fi
\noindent\textit{Source reference.} Original lemma 0EEW: focused-seeds/AGP/sources/stacks-snapshot-20261011/tags.txt, source lines 18041--18041. Stable Stacks tag for an expressly allowed foundational or notational input; the counted construction remains the two concrete core obligations. The retained local core is not replaced by this pointer. See standard\_prerequisites.\par\medskip
\ifdefined\hypertarget\hypertarget{sourcepointer-18}{}\fi
\noindent\textit{Source reference.} Original lemma 0EEX: focused-seeds/AGP/sources/stacks-snapshot-20261011/tags.txt, source lines 18042--18042. Stable Stacks tag for an expressly allowed foundational or notational input; the counted construction remains the two concrete core obligations. The retained local core is not replaced by this pointer. See standard\_prerequisites.\par\medskip
\ifdefined\hypertarget\hypertarget{sourcepointer-19}{}\fi
\noindent\textit{Source reference.} Original lemma 0BK4: focused-seeds/AGP/sources/stacks-snapshot-20261011/tags.txt, source lines 14547--14547. Stable Stacks tag for an expressly allowed foundational or notational input; the counted construction remains the two concrete core obligations. The retained local core is not replaced by this pointer. See standard\_prerequisites.\par\medskip
\ifdefined\hypertarget\hypertarget{sourcepointer-20}{}\fi
\noindent\textit{Source reference.} Original section 0EF4: focused-seeds/AGP/sources/stacks-snapshot-20261011/tags.txt, source lines 18049--18049. Stable Stacks tag for an expressly allowed foundational or notational input; the counted construction remains the two concrete core obligations. The retained local core is not replaced by this pointer. See standard\_prerequisites.\par\medskip
\begin{lemma}
\label{lemma-algebraize-local-cohomology-general}
Let $I, J$ be ideals of a Noetherian ring $A$.
Let $M$ be a finite $A$-module.
Let $s$ and $d$ be integers. With $T$ as in
(\ref{equation-associated-subset}) assume
\begin{enumerate}
\item $A$ is $I$-adically complete and has a dualizing complex,
\item if $\mathfrak p \in V(I)$ no condition,
\item $\text{cd}(A, I) \leq d$,
\item if $\mathfrak p \not \in V(I)$, $\mathfrak p \not \in T$ then
$$
\text{depth}_{A_\mathfrak p}(M_\mathfrak p) \geq s
\quad\text{or}\quad
\text{depth}_{A_\mathfrak p}(M_\mathfrak p) +
\dim((A/\mathfrak p)_\mathfrak q) > d + s
$$
for all $\mathfrak q \in V(\mathfrak p) \cap V(J) \cap V(I)$,
\item if $\mathfrak p \not \in V(I)$, $\mathfrak p \not \in T$,
$V(\mathfrak p) \cap V(J) \cap V(I) \not = \emptyset$, and
$\text{depth}(M_\mathfrak p) < s$, then one
of the following holds\footnote{Our method
forces this additional condition. We will return to this
(insert future reference).}:
\begin{enumerate}
\item $\dim(\text{Supp}(M_\mathfrak p)) < s + 2$\footnote{For example
if $M$ satisfies Serre's condition $(S_s)$
on the complement of $V(I) \cup T$.}, or
\item  $\delta(\mathfrak p) > d + \delta_{max} - 1$
where $\delta$ is a dimension function and $\delta_{max}$
is the maximum of $\delta$ on $V(J) \cap V(I)$, or
\item $\text{depth}_{A_\mathfrak p}(M_\mathfrak p) +
\dim((A/\mathfrak p)_\mathfrak q) > d + s + \delta_{max} - \delta_{min} - 2$
for all $\mathfrak q \in V(\mathfrak p) \cap V(J) \cap V(I)$.
\end{enumerate}
\end{enumerate}
Then there exists an ideal $J_0 \subset J$ with
$V(J_0) \cap V(I) = V(J) \cap V(I)$
such that for any $J' \subset J_0$ with
$V(J') \cap V(I) = V(J) \cap V(I)$ the map
$$
R\Gamma_{J'}(M) \longrightarrow R\Gamma_J(M)^\wedge
$$
induces an isomorphism on cohomology in degrees $\leq s$.
Here ${}^\wedge$ denotes derived $I$-adic completion.
\end{lemma}

\section{Algebraization of local cohomology, I}
\label{section-algebraization-sections-general}

\noindent
Let $A$ be a Noetherian ring and let $I$ and $J$ be two ideals of $A$.
Let $M$ be a finite $A$-module. In this section we study the
cohomology groups of the object
$$
R\Gamma_J(M)^\wedge
\quad\text{of}\quad
D(A)
$$
where ${}^\wedge$ denotes derived $I$-adic completion. Observe that in
Dualizing Complexes, Lemma \ref{dualizing-lemma-completion-local-H0}
we have shown, if $A$ is complete with respect to $I$,
that there is an isomorphism
$$
\colim H^0_Z(M) \longrightarrow H^0(R\Gamma_J(M)^\wedge)
$$
where the (directed) colimit is over the closed subsets $Z = V(J')$
with $J' \subset J$ and $V(J') \cap V(I) = V(J) \cap V(I)$.
The union of these closed subsets is
\begin{equation}
\label{equation-associated-subset}
T = \{\mathfrak p \in \Spec(A) :
V(\mathfrak p) \cap V(I) \subset V(J) \cap V(I)\}
\end{equation}
This is a subset of $\Spec(A)$ stable under specialization.
The result above becomes the statement that
$$
H^0_T(M) \longrightarrow H^0(R\Gamma_J(M)^\wedge)
$$
is an isomorphism provided $A$ is complete with respect to $I$, see
Local Cohomology, Lemma \ref{local-cohomology-lemma-adjoint-ext} and
Remark \ref{local-cohomology-remark-upshot}.
Our method to extend this isomorphism to higher cohomology groups
rests on the following lemma.

\begin{lemma}
\label{lemma-kill-completion-general}
Let $I, J$ be ideals of a Noetherian ring $A$.
Let $M$ be a finite $A$-module. Let $\mathfrak p \subset A$ be a prime.
Let $s$ and $d$ be integers. Assume
\begin{enumerate}
\item $A$ has a dualizing complex,
\item $\mathfrak p \not \in V(J) \cap V(I)$,
\item $\text{cd}(A, I) \leq d$, and
\item for all primes $\mathfrak p' \subset \mathfrak p$
we have
$\text{depth}_{A_{\mathfrak p'}}(M_{\mathfrak p'}) +
\dim((A/\mathfrak p')_\mathfrak q) > d + s$
for all $\mathfrak q \in V(\mathfrak p') \cap V(J) \cap V(I)$.
\end{enumerate}
Then there exists an $f \in A$, $f \not \in \mathfrak p$ which annihilates
$H^i(R\Gamma_J(M)^\wedge)$ for $i \leq s$ where ${}^\wedge$
indicates $I$-adic completion.
\end{lemma}

\begin{proof}
We will use that $R\Gamma_J = R\Gamma_{V(J)}$ and similarly for
$I + J$, see
Dualizing Complexes, Lemma \ref{dualizing-lemma-local-cohomology-noetherian}.
Observe that
$R\Gamma_J(M)^\wedge = R\Gamma_I(R\Gamma_J(M))^\wedge =
R\Gamma_{I + J}(M)^\wedge$, see
Dualizing Complexes, Lemmas
\ref{dualizing-lemma-complete-and-local} and
\ref{dualizing-lemma-local-cohomology-ss}.
Thus we may replace $J$ by $I + J$ and assume $I \subset J$
and $\mathfrak p \not \in V(J)$.
Recall that
$$
R\Gamma_J(M)^\wedge = R\Hom_A(R\Gamma_I(A), R\Gamma_J(M))
$$
by the description of derived completion in
More on Algebra, Lemma \ref{more-algebra-lemma-derived-completion}
combined with the description of local cohomology in
Dualizing Complexes, Lemma
\ref{dualizing-lemma-compute-local-cohomology-noetherian}.
Assumption (3) means that $R\Gamma_I(A)$ has nonzero cohomology
only in degrees $\leq d$. Using the canonical truncations of
$R\Gamma_I(A)$ we find it suffices to show that
$$
\text{Ext}^i(N, R\Gamma_J(M))
$$
is annihilated by an $f \in A$, $f \not \in \mathfrak p$ for
$i \leq s + d$ and any $A$-module $N$.
In turn using the canonical truncations for $R\Gamma_J(M)$
we see that it suffices to show
$H^i_J(M)$ is annihilated by an $f \in A$, $f \not \in \mathfrak p$
for $i \leq s + d$.
This follows from Local Cohomology, Lemma
\ref{local-cohomology-lemma-kill-local-cohomology-at-prime}.
\end{proof}

\begin{lemma}
\label{lemma-kill-colimit-weak-general}
Let $I, J$ be ideals of a Noetherian ring. Let $M$ be a finite $A$-module.
Let $s$ and $d$ be integers. With $T$ as in
(\ref{equation-associated-subset}) assume
\begin{enumerate}
\item $A$ has a dualizing complex,
\item if $\mathfrak p \in V(I)$, then no condition,
\item if $\mathfrak p \not \in V(I)$, $\mathfrak p \in T$, then
$\dim((A/\mathfrak p)_\mathfrak q) \leq d$ for some
$\mathfrak q \in V(\mathfrak p) \cap V(J) \cap V(I)$,
\item if $\mathfrak p \not \in V(I)$, $\mathfrak p \not \in T$, then
$$
\text{depth}_{A_\mathfrak p}(M_\mathfrak p) \geq s
\quad\text{or}\quad
\text{depth}_{A_\mathfrak p}(M_\mathfrak p) +
\dim((A/\mathfrak p)_\mathfrak q) > d + s
$$
for all $\mathfrak q \in V(\mathfrak p) \cap V(J) \cap V(I)$.
\end{enumerate}
Then there exists an ideal $J_0 \subset J$ with
$V(J_0) \cap V(I) = V(J) \cap V(I)$ such that for any $J' \subset J_0$ with
$V(J') \cap V(I) = V(J) \cap V(I)$ the map
$$
R\Gamma_{J'}(M) \longrightarrow R\Gamma_{J_0}(M)
$$
induces an isomorphism in cohomology in degrees $\leq s$
and moreover these modules are annihilated by a power of $J_0I$.
\end{lemma}

\begin{proof}
Let us consider the set
$$
B = \{\mathfrak p \not \in V(I),\ \mathfrak p \in T,\text{ and }
\text{depth}(M_\mathfrak p) \leq s\}
$$
Choose $J_0 \subset J$ such that $V(J_0)$ is the closure of $B \cup V(J)$.

\medskip\noindent
Claim I: $V(J_0) \cap V(I) = V(J) \cap V(I)$.

\medskip\noindent
Proof of Claim I. The inclusion $\supset$ holds by construction.
Let $\mathfrak p$ be a minimal prime of $V(J_0)$.
If $\mathfrak p \in B \cup V(J)$, then either $\mathfrak p \in T$
or $\mathfrak p \in V(J)$ and in both cases
$V(\mathfrak p) \cap V(I) \subset V(J) \cap V(I)$ as desired.
If $\mathfrak p \not \in B \cup V(J)$, then
$V(\mathfrak p) \cap B$ is dense, hence infinite, and we conclude that
$\text{depth}(M_\mathfrak p) < s$ by
Local Cohomology, Lemma \ref{local-cohomology-lemma-depth-function}.
In fact, let
$V(\mathfrak p) \cap B = \{\mathfrak p_\lambda\}_{\lambda \in \Lambda}$.
Pick $\mathfrak q_\lambda \in V(\mathfrak p_\lambda) \cap V(J) \cap V(I)$
as in (3).
Let $\delta : \Spec(A) \to \mathbf{Z}$ be the dimension function
associated to a dualizing complex $\omega_A^\bullet$ for $A$.
Since $\Lambda$ is infinite and $\delta$ is bounded,
there exists an infinite subset $\Lambda' \subset \Lambda$ on which
$\delta(\mathfrak q_\lambda)$ is constant. For
$\lambda \in \Lambda'$ we have
$$
\text{depth}(M_{\mathfrak p_\lambda}) +
\delta(\mathfrak p_\lambda) - \delta(\mathfrak q_\lambda) =
\text{depth}(M_{\mathfrak p_\lambda}) +
\dim((A/\mathfrak p_\lambda)_{\mathfrak q_\lambda})
\leq d + s
$$
by (3) and the definition of $B$. By the semi-continuity of
the function $\text{depth} + \delta$ proved in
Duality for Schemes, Lemma \ref{duality-lemma-sitting-in-degrees}
we conclude that
$$
\text{depth}(M_\mathfrak p) +
\dim((A/\mathfrak p)_{\mathfrak q_\lambda}) =
\text{depth}(M_\mathfrak p) + \delta(\mathfrak p) - \delta(\mathfrak q_\lambda)
\leq d + s
$$
Since also $\mathfrak p \not \in V(I)$ we read off from (4) that
$\mathfrak p \in T$, i.e.,
$V(\mathfrak p) \cap V(I) \subset V(J) \cap V(I)$. This finishes the
proof of Claim I.

\medskip\noindent
Claim II: $H^i_{J_0}(M) \to H^i_J(M)$ is an isomorphism for $i \leq s$
and $J' \subset J_0$ with $V(J') \cap V(I) = V(J) \cap V(I)$.

\medskip\noindent
Proof of claim II. Choose $\mathfrak p \in V(J')$ not in $V(J_0)$.
It suffices to show that $H^i_{\mathfrak pA_\mathfrak p}(M_\mathfrak p) = 0$
for $i \leq s$, see
Local Cohomology, Lemma \ref{local-cohomology-lemma-isomorphism}.
Observe that $\mathfrak p \in T$. Hence since $\mathfrak p$ is not in $B$
we see that $\text{depth}(M_\mathfrak p) > s$ and the groups vanish by
Dualizing Complexes, Lemma \ref{dualizing-lemma-depth}.

\medskip\noindent
Claim III. The final statement of the lemma is true.

\medskip\noindent
By Claim II for $i \leq s$ we have
$$
H^i_T(M) = H^i_{J_0}(M) = H^i_{J'}(M)
$$
for all ideals $J' \subset J_0$ with $V(J')  \cap V(I) = V(J) \cap V(I)$.
See Local Cohomology, Lemma \ref{local-cohomology-lemma-adjoint-ext}.
Let us check the hypotheses of Local Cohomology,
Proposition \ref{local-cohomology-proposition-annihilator}
for the subsets $T \subset T \cup V(I)$, the module $M$, and the integer $s$.
We have to show that given $\mathfrak p \subset \mathfrak q$
with $\mathfrak p \not \in T \cup V(I)$ and $\mathfrak q \in T$
we have
$$
\text{depth}_{A_\mathfrak p}(M_\mathfrak p) +
\dim((A/\mathfrak p)_\mathfrak q) > s
$$
If $\text{depth}(M_\mathfrak p) \geq s$, then this is true because
the dimension of $(A/\mathfrak p)_\mathfrak q$ is at least $1$.
Thus we may assume $\text{depth}(M_\mathfrak p) < s$.
If $\mathfrak q \in V(I)$, then $\mathfrak q \in V(J) \cap V(I)$
and the inequality holds by (4). If $\mathfrak q \not \in V(I)$,
then we can use (3) to pick
$\mathfrak q' \in V(\mathfrak q) \cap V(J) \cap V(I)$ with
$\dim((A/\mathfrak q)_{\mathfrak q'}) \leq d$.
Then assumption (4) gives
$$
\text{depth}_{A_\mathfrak p}(M_\mathfrak p) +
\dim((A/\mathfrak p)_{\mathfrak q'}) > s + d
$$
Since $A$ is catenary this implies the inequality we want.
Applying Local Cohomology,
Proposition \ref{local-cohomology-proposition-annihilator} we
find $J'' \subset A$ with $V(J'') \subset T \cup V(I)$
such that $J''$ annihilates $H^i_T(M)$ for $i \leq s$.
Then we can write $V(J'') \cup V(J_0) \cup V(I) = V(J'I)$
for some $J' \subset J_0$ with $V(J') \cap V(I) = V(J) \cap V(I)$.
Replacing $J_0$ by $J'$ the proof is complete.
\end{proof}

\begin{lemma}
\label{lemma-kill-colimit-general}
In Lemma \ref{lemma-kill-colimit-weak-general} if instead of the empty
condition (2) we assume
\begin{enumerate}
\item[(2')] if $\mathfrak p \in V(I)$, $\mathfrak p \not \in V(J) \cap V(I)$,
then
$\text{depth}_{A_\mathfrak p}(M_\mathfrak p) +
\dim((A/\mathfrak p)_\mathfrak q) > s$
for all $\mathfrak q \in V(\mathfrak p) \cap V(J) \cap V(I)$,
\end{enumerate}
then the conditions also imply that $H^i_{J_0}(M)$ is a finite
$A$-module for $i \leq s$.
\end{lemma}

\begin{proof}
Recall that $H^i_{J_0}(M) = H^i_T(M)$, see proof of
Lemma \ref{lemma-kill-colimit-weak-general}. Thus it suffices to
check that for $\mathfrak p \not \in T$ and $\mathfrak q \in T$
with $\mathfrak p \subset \mathfrak q$ we have
$\text{depth}_{A_\mathfrak p}(M_\mathfrak p) +
\dim((A/\mathfrak p)_\mathfrak q) > s$, see Local Cohomology,
Proposition \ref{local-cohomology-proposition-finiteness}.
Condition (2') tells us this is true for $\mathfrak p \in V(I)$.
Since we know $H^i_T(M)$ is annihilated by a power of $IJ_0$
we know the condition holds if $\mathfrak p \not \in V(IJ_0)$
by Local Cohomology, Proposition \ref{local-cohomology-proposition-annihilator}.
This covers all cases and the proof is complete.
\end{proof}

\begin{lemma}
\label{lemma-kill-colimit-support-general}
If in Lemma \ref{lemma-kill-colimit-weak-general} we additionally assume
\begin{enumerate}
\item[(6)] if $\mathfrak p \not \in V(I)$, $\mathfrak p \in T$, then
$\text{depth}_{A_\mathfrak p}(M_\mathfrak p) > s$,
\end{enumerate}
then $H^i_{J_0}(M) = H^i_J(M) = H^i_{J + I}(M)$ for $i \leq s$ and these
modules are annihilated by a power of $I$.
\end{lemma}

\begin{proof}
Choose $\mathfrak p \in V(J)$ or $\mathfrak p \in V(J_0)$ but
$\mathfrak p \not \in V(J + I) = V(J_0 + I)$.
It suffices to show that $H^i_{\mathfrak pA_\mathfrak p}(M_\mathfrak p) = 0$
for $i \leq s$, see
Local Cohomology, Lemma \ref{local-cohomology-lemma-isomorphism}.
These groups vanish by condition (6) and
Dualizing Complexes, Lemma \ref{dualizing-lemma-depth}.
The final statement follows from
Local Cohomology, Proposition \ref{local-cohomology-proposition-annihilator}.
\end{proof}

\begin{proof}
For an ideal $\mathfrak a \subset A$ we have
$R\Gamma_\mathfrak a = R\Gamma_{V(\mathfrak a)}$, see
Dualizing Complexes, Lemma \ref{dualizing-lemma-local-cohomology-noetherian}.
Next, we observe that
$$
R\Gamma_J(M)^\wedge =
R\Gamma_I(R\Gamma_J(M))^\wedge =
R\Gamma_{I + J}(M)^\wedge =
R\Gamma_{I + J'}(M)^\wedge =
R\Gamma_I(R\Gamma_{J'}(M))^\wedge =
R\Gamma_{J'}(M)^\wedge
$$
by Dualizing Complexes, Lemmas \ref{dualizing-lemma-local-cohomology-ss} and
\ref{dualizing-lemma-complete-and-local}.
This explains how we define the arrow in the statement of the lemma.

\medskip\noindent
We claim that the hypotheses of Lemma \ref{lemma-kill-colimit-weak-general}
are implied by our current hypotheses on $M$.
The only thing to verify is hypothesis (3).
Thus let $\mathfrak p \not \in V(I)$, $\mathfrak p \in T$.
Then $V(\mathfrak p) \cap V(I)$ is nonempty as $I$ is
contained in the Jacobson radical of $A$
(Algebra, Lemma \ref{algebra-lemma-radical-completion}).
Since $\mathfrak p \in T$ we have
$V(\mathfrak p) \cap V(I) = V(\mathfrak p) \cap V(J) \cap V(I)$.
Let $\mathfrak q \in V(\mathfrak p) \cap V(I)$ be the
generic point of an irreducible component.
We have $\text{cd}(A_\mathfrak q, I_\mathfrak q) \leq d$
by Local Cohomology, Lemma \ref{local-cohomology-lemma-cd-local}.
We have $V(\mathfrak pA_\mathfrak q) \cap V(I_\mathfrak q) =
\{\mathfrak qA_\mathfrak q\}$ by our choice of $\mathfrak q$
and we conclude $\dim((A/\mathfrak p)_\mathfrak q) \leq d$
by Local Cohomology, Lemma \ref{local-cohomology-lemma-cd-bound-dim-local}.

\medskip\noindent
Observe that the lemma holds for $s < 0$. This is not a trivial case because
it is not a priori clear that $H^i(R\Gamma_J(M)^\wedge)$
is zero for $i < 0$. However, this vanishing was established in
Dualizing Complexes, Lemma \ref{dualizing-lemma-completion-local}.
We will prove the lemma by induction for $s \geq 0$.

\medskip\noindent
The lemma for $s = 0$ follows immediately from
the conclusion of Lemma \ref{lemma-kill-colimit-weak-general}
and Dualizing Complexes, Lemma \ref{dualizing-lemma-completion-local-H0}.

\medskip\noindent
Assume $s > 0$ and the lemma has been shown for smaller values of $s$.
Let $M' \subset M$ be the maximal submodule whose support is contained
in $V(I) \cup T$. Then $M'$ is a finite $A$-module whose support
is contained in $V(J') \cup V(I)$ for some ideal $J' \subset J$
with $V(J') \cap V(I) = V(J) \cap V(I)$.
We claim that
$$
R\Gamma_{J'}(M') \to R\Gamma_J(M')^\wedge
$$
is an isomorphism for any choice of $J'$.
Namely, we can choose a short exact sequence
$0 \to M_1 \oplus M_2 \to M' \to N \to 0$ with
$M_1$ annihilated by a power of $J'$, with $M_2$ annihilated
by a power of $I$, and with $N$ annihilated by a power of $I + J'$.
Thus it suffices to show that the claim holds for $M_1$, $M_2$, and $N$.
In the case of $M_1$ we see that $R\Gamma_{J'}(M_1) = M_1$ and
since $M_1$ is a finite $A$-module and $I$-adically complete
we have $M_1^\wedge = M_1$. This proves the claim for $M_1$
by the initial remarks of the proof. In the case of $M_2$ we see that
$H^i_J(M_2) = H^i_{I + J}(M) = H^i_{I + J'}(M) = H^i_{J'}(M_2)$
are annihilated by a power of $I$ and hence derived complete.
Thus the claim in this case also. For $N$ we can use either of
the arguments just given. Considering the short exact sequence
$0 \to M' \to M \to M/M' \to 0$
we see that it suffices to prove the lemma for $M/M'$.
Thus we may assume $\text{Ass}(M) \cap (V(I) \cup T) = \emptyset$.

\medskip\noindent
Let $\mathfrak p \in \text{Ass}(M)$ be such that
$V(\mathfrak p) \cap V(J) \cap V(I) = \emptyset$.
Since $I$ is contained in the Jacobson radical of $A$ this implies
that $V(\mathfrak p) \cap V(J') = \emptyset$ for any
$J' \subset J$ with $V(J') \cap V(I) = V(J) \cap V(I)$.
Thus setting $N = H^0_\mathfrak p(M)$ we see that
$R\Gamma_J(N) = R\Gamma_{J'}(N) = 0$ for all
$J' \subset J$ with $V(J') \cap V(I) = V(J) \cap V(I)$.
In particular $R\Gamma_J(N)^\wedge = 0$.
Thus we may replace $M$ by $M/N$ as this changes the
structure of $M$ only in primes which do not play
a role in conditions (4) or (5). Repeating we may assume that
$V(\mathfrak p) \cap V(J) \cap V(I) \not = \emptyset$
for all $\mathfrak p \in \text{Ass}(M)$.

\medskip\noindent
Assume $\text{Ass}(M) \cap (V(I) \cup T) = \emptyset$ and that
$V(\mathfrak p) \cap V(J) \cap V(I) \not = \emptyset$
for all $\mathfrak p \in \text{Ass}(M)$.
Let $\mathfrak p \in \text{Ass}(M)$. We want to show that we may apply
Lemma \ref{lemma-kill-completion-general}.
It is in the verification of this that we will use the supplemental
condition (5). Choose $\mathfrak p' \subset \mathfrak p$
and $\mathfrak q' \subset V(\mathfrak p) \cap V(J) \cap V(I)$.
\begin{enumerate}
\item If $M_{\mathfrak p'} = 0$, then
$\text{depth}(M_{\mathfrak p'}) = \infty$ and
$\text{depth}(M_{\mathfrak p'}) +
\dim((A/\mathfrak p')_{\mathfrak q'}) > d + s$.
\item If $\text{depth}(M_{\mathfrak p'}) < s$, then
$\text{depth}(M_{\mathfrak p'}) +
\dim((A/\mathfrak p')_{\mathfrak q'}) > d + s$ by (4).
\end{enumerate}
In the remaining cases we have $M_{\mathfrak p'} \not = 0$ and
$\text{depth}(M_{\mathfrak p'}) \geq s$. In particular, we see that
$\mathfrak p'$ is in the support of $M$ and we can choose
$\mathfrak p'' \subset \mathfrak p'$ with $\mathfrak p'' \in \text{Ass}(M)$.
\begin{enumerate}
\item[(a)] Observe that
$\dim((A/\mathfrak p'')_{\mathfrak p'}) \geq \text{depth}(M_{\mathfrak p'})$
by Algebra, Lemma \ref{algebra-lemma-depth-dim-associated-primes}.
If equality holds, then we have
$$
\text{depth}(M_{\mathfrak p'}) + \dim((A/\mathfrak p')_{\mathfrak q'}) =
\text{depth}(M_{\mathfrak p''}) + \dim((A/\mathfrak p'')_{\mathfrak q'})
> s + d
$$
by (4) applied to $\mathfrak p''$ and we are done. This means we are
only in trouble if
$\dim((A/\mathfrak p'')_{\mathfrak p'}) > \text{depth}(M_{\mathfrak p'})$.
This implies that $\dim(M_\mathfrak p) \geq s + 2$.
Thus if (5)(a) holds, then this does not occur.
\item[(b)] If (5)(b) holds, then we get
$$
\text{depth}(M_{\mathfrak p'}) + \dim((A/\mathfrak p')_{\mathfrak q'})
\geq s + \delta(\mathfrak p') - \delta(\mathfrak q')
\geq s + 1 + \delta(\mathfrak p) - \delta_{max}
> s + d
$$
as desired.
\item[(c)] If (5)(c) holds, then we get
\begin{align*}
\text{depth}(M_{\mathfrak p'}) + \dim((A/\mathfrak p')_{\mathfrak q'})
& \geq
s + \delta(\mathfrak p') - \delta(\mathfrak q') \\
& \geq
s + 1 + \delta(\mathfrak p) - \delta(\mathfrak q') \\
& =
s + 1 + \delta(\mathfrak p) - \delta(\mathfrak q) +
\delta(\mathfrak q) - \delta(\mathfrak q') \\
& >
s + 1 + (s + d + \delta_{max} - \delta_{min} - 2) +
\delta(\mathfrak q) - \delta(\mathfrak q') \\
& \geq 
2s + d - 1 \geq s + d
\end{align*}
as desired. Observe that this argument works because
we know that a prime $\mathfrak q \in V(\mathfrak p) \cap V(J) \cap V(I)$
exists.
\end{enumerate}
Now we are ready to do the induction step.

\medskip\noindent
Choose an ideal $J_0$ as in Lemma \ref{lemma-kill-colimit-weak-general}
and an integer $t > 0$ such that $(J_0I)^t$ annihilates $H^s_J(M)$.
The assumptions of Lemma \ref{lemma-kill-completion-general}
are satisfied for every $\mathfrak p \in \text{Ass}(M)$
(see previous paragraph).
Thus the annihilator $\mathfrak a \subset A$ of
$H^s(R\Gamma_J(M)^\wedge)$
is not contained in $\mathfrak p$ for $\mathfrak p \in \text{Ass}(M)$.
Thus we can find an $f \in \mathfrak a(J_0I)^t$
not in any associated prime of $M$ which is an annihilator
of both $H^s(R\Gamma_J(M)^\wedge)$ and $H^s_J(M)$.
Then $f$ is a nonzerodivisor on $M$ and we can consider the
short exact sequence
$$
0 \to M \xrightarrow{f} M \to M/fM \to 0
$$
Our choice of $f$ shows that we obtain
$$
\xymatrix{
H^{s - 1}_{J'}(M) \ar[d] \ar[r] &
H^{s - 1}_{J'}(M/fM) \ar[d] \ar[r] &
H^s_{J'}(M) \ar[d] \ar[r] & 0 \\
H^{s - 1}(R\Gamma_J(M)^\wedge) \ar[r] &
H^{s - 1}(R\Gamma_J(M/fM)^\wedge) \ar[r] &
H^s(R\Gamma_J(M)^\wedge) \ar[r] & 0
}
$$
for any $J' \subset J_0$ with $V(J') \cap V(I) = V(J) \cap V(I)$.
Thus if we choose $J'$ such that it works for
$M$ and $M/fM$ and $s - 1$ (possible by induction hypothesis --
see next paragraph), then we conclude that the lemma is true.

\medskip\noindent
To finish the proof we have to show that the module
$M/fM$ satisfies the hypotheses (4) and (5) for $s - 1$.
Thus we let $\mathfrak p$ be a prime in the support
of $M/fM$ with $\text{depth}((M/fM)_\mathfrak p) < s - 1$
and with $V(\mathfrak p) \cap V(J) \cap V(I)$ nonempty.
Then $\dim(M_\mathfrak p) = \dim((M/fM)_\mathfrak p) + 1$
and $\text{depth}(M_\mathfrak p) = \text{depth}((M/fM)_\mathfrak p) + 1$.
In particular, we know (4) and (5) hold for $\mathfrak p$ and $M$
with the original value $s$.
The desired inequalities then follow by inspection.
\end{proof}

\bibliographystyle{plain}
\bibliography{my}
\end{document}
